\chapter[RAPPORTS MÉTÉO]{Rapports météo}\label{rapports-muxe9tuxe9o}

APRS supporte Peet Brothers, Ultimeter et Davis. Idéal pour Skywarn.

\section{Trois types}\label{trois-types}

\begin{itemize}
\tightlist
\item
  \textbf{Raw Weather Report}
\item
  \textbf{Positionless Weather Report}
\item
  \textbf{Complete Weather Report} (avec position et timestamp)
\end{itemize}

\section{DTIs des Weather Reports}\label{dtis-des-weather-reports}

Pour données brutes :

\begin{itemize}
\tightlist
\item
  \texttt{!} : Ultimeter 2000
\item
  \texttt{\#} : Peet Bros U-II
\item
  \texttt{\$} : Ultimeter 2000
\item
  \texttt{*} : Peet Bros U-II
\item
  \texttt{\_} : Positionless weather data
\end{itemize}

Après post-traitement (insertion position), on peut utiliser \texttt{!},
\texttt{=}, \texttt{/}, \texttt{@} classiques. Le rapport est alors
identifié par le symbole weather \texttt{/\_} ou
\texttt{\textbackslash{}\_} dans les APRS Data.

\section{Raw Weather Report}\label{raw-weather-report}

\begin{verbatim}
! ou # ou $ ou * | Raw Weather Data (n octets)
\end{verbatim}

Exemples : -
\texttt{!!006B005803500000-\/-\/-\/-03E9-\/-\/-\/-\/-\/-\/-\/-002105140000005D}
(Ultimeter 2000) - \texttt{\#50B7500820082} (Peet Bros U-II) -
\texttt{\$ULTW0031003702CE0069-\/-\/-\/-000086A00001-\/-\/-\/-011901CC00000005}
(Ultimeter 2000) - \texttt{*7007600000000} (Peet Bros U-II)

\section{Positionless Weather
Report}\label{positionless-weather-report}

\begin{center}
\begin{tabular}{|c|c|c|c|c|c|}
\hline
\textbf{Champ} & \lit{\_} & Time MDHM & Positionless Weather Data & APRS Software (S) & WX Unit (uuuu) \\
\hline
\textbf{Octets} & 1 & 8 & n & 1 & 2-4 \\
\hline
\end{tabular}
\end{center}

Ex : \texttt{\_10090556c220s004g005t077r000p000P000h50b09900wRSW}.

Codes logiciel : \texttt{d} APRSdos, \texttt{M} MacAPRS, \texttt{P}
pocketAPRS, \texttt{S} APRS+SA, \texttt{W} WinAPRS, \texttt{X} X-APRS.

Codes unité (2-4 chars) : \texttt{Dvs}, \texttt{HKT}, \texttt{PIC},
\texttt{RSW}, \texttt{U-II}, \texttt{U2R}, \texttt{U2k}, \texttt{U2kr},
\texttt{U5}, \texttt{Upkm}. Codes libres possibles.

\section{Format des données météo
positionless}\label{format-des-donnuxe9es-muxe9tuxe9o-positionless}

\begin{center}
\begin{tabular}{|c|c|c|c|c|c|c|c|c|}
\hline
\textbf{Champ} & Wind dir & Wind speed & Gust & Temp & Rain hr & Rain 24h & Rain midnight & Humidity \\
\textbf{Code} & \texttt{cccc} & \texttt{ssss} & \texttt{gggg} & \texttt{tttt} & \texttt{rrrr} & \texttt{pppp} & \texttt{PPPP} & \texttt{hhh} \\
\textbf{Octets} & 4 & 4 & 4 & 4 & 4 & 4 & 4 & 3 \\
\hline
\end{tabular}
\end{center}
\noindent + Barometric Pressure (\texttt{bbbbbb}, 5 octets)

\begin{itemize}
\tightlist
\item
  \texttt{cccc} = direction vent en degrés.
\item
  \texttt{ssss} = vitesse soutenue 1 min en mph.
\item
  \texttt{gggg} = pic de rafale (5 min) en mph.
\item
  \texttt{tttt} = température en °F. En dessous de zéro : \texttt{-01} à
  \texttt{-99}.
\item
  \texttt{rrrr}, \texttt{pppp}, \texttt{PPPP} = pluie en 1/100 pouce
  (dernière heure, 24 h, depuis minuit).
\item
  \texttt{hhh} = humidité en \% (00 = 100 \%).
\item
  \texttt{bbbbbb} = pression barométrique en 1/10 millibars.
\end{itemize}

Paramètres additionnels : - \texttt{L} : luminosité W/m² (≤ 999). -
\texttt{l} : luminosité ≥ 1000 (remplace un champ de pluie). -
\texttt{s} : neige 24 h (pouces). - \texttt{\#} : compteur brut de
pluie.

Note : au minimum MDHM + wind dir + speed + gust + temp doivent être
présents (même invalides). Reste dans un ordre libre ou absent. Valeur
inconnue = points ou espaces (\texttt{c...s...g...} ou
\texttt{cVVVsVVVgVVV}).

Ex pluie seule : \texttt{\_10090556c...s...g...t...P012Jim}.

\section{Localisation d'un
raw/positionless}\label{localisation-dun-rawpositionless}

Le rapport n'a pas de position. La station doit occasionnellement
émettre un paquet séparé avec sa position.

\section{Symbole d'un
raw/positionless}\label{symbole-dun-rawpositionless}

Le raw/positionless n'a pas de Symbol Code dans l'info. Spécifier via
une adresse de destination APRS générique (ex \texttt{GPSHW} ou
\texttt{GPSE63}) ou, sur ballon, via SSID \texttt{-11} sur l'adresse
source. Voir chapitre 20.

\section{Complete Weather Report}\label{complete-weather-report}

Symbol Code toujours \texttt{\_}. Le 7-byte Wind Direction/Speed
extension remplace \texttt{cccc/ssss} du positionless.

\textbf{Avec lat/long, sans timestamp} :

\begin{center}
\begin{tabular}{|c|c|c|c|c|c|c|c|c|c|}
\hline
\textbf{Champ} & \lit{!} ou \lit{=} & Lat & Sym Tbl ID & Long & \lit{\_} & Wind dir/speed & Weather Data & S & uuuu \\
\hline
\textbf{Octets} & 1 & 8 & 1 & 9 & 1 & 7 & n & 1 & 2-4 \\
\hline
\end{tabular}
\end{center}

Ex :
\texttt{!4903.50N/07201.75W\_220/004g005t077r000p000P000h50b09900wRSW}.

\textbf{Avec lat/long et timestamp} :

\begin{center}
\begin{tabular}{|c|c|c|c|c|c|c|c|c|c|c|}
\hline
\textbf{Champ} & \lit{/} ou \lit{@} & Time & Lat & Sym Tbl ID & Long & \lit{\_} & Wind & Weather Data & S & uuuu \\
\hline
\textbf{Octets} & 1 & 7 & 8 & 1 & 9 & 1 & 7 & n & 1 & 2-4 \\
\hline
\end{tabular}
\end{center}

Ex :
\texttt{@092345z4903.50N/07201.75W\_220/004g005t-07r000p000P000h50b09900wRSW}.

\textbf{Avec position compressée} : variantes équivalentes.

\textbf{Weather via Object Report} :

\begin{center}
\begin{tabular}{|c|c|c|c|c|c|c|c|c|c|}
\hline
\textbf{Champ} & \lit{;} & Object Name & \lit{*} & Time & Lat & Sym Tbl ID & Long & \lit{\_} & Wind + Weather Data \\
\hline
\textbf{Octets} & 1 & 9 & 1 & 7 & 8 & 1 & 9 & 1 & 7 + n \\
\hline
\end{tabular}
\end{center}

Ex :
\texttt{;BRENDA\ \ \ *4903.50N/07201.75W\_220/004g005t077r000p000P000h50b09900wRSW}.

\section{Storm Data}\label{storm-data}

Format des données de tempête :

\begin{center}
\begin{tabular}{|c|c|c|c|c|c|c|c|}
\hline
\textbf{Champ} & Direction & Storm Type & Sustained wind & Peak gust & Central pressure & Hurricane radius & TS radius \\
\textbf{Séparateur} &  & \texttt{/ST} & \texttt{/www} & \texttt{\^{}GGG} & \texttt{/pppp} & \texttt{>RRR} & \texttt{\&rrr} \\
\textbf{Octets} & 3 & 3 & 4 & 4 & 5 & 4 & 4 \\
\hline
\end{tabular}
\end{center}
\noindent + Whole Gale radius (\texttt{\%ggg}, 4 octets, optionnel)

\begin{itemize}
\tightlist
\item
  \texttt{ST} : TS = Tropical Storm, HC = Hurricane, TD = Tropical
  Depression.
\item
  \texttt{www} : vent soutenu en knots.
\item
  \texttt{GGG} : pic de rafale en knots.
\item
  \texttt{pppp} : pression centrale en millibars.
\item
  \texttt{RRR} : rayon vents ouragan en milles nautiques.
\item
  \texttt{rrr} : rayon vents tempête tropicale en milles nautiques.
\item
  \texttt{ggg} : rayon vents « whole gale » (50 knots) en milles
  nautiques. Optionnel.
\end{itemize}

Symboles : - \texttt{\textbackslash{}@} : Hurricane/Tropical Storm
(position actuelle). - \texttt{/@} : position future prédite.

Ex progression Hurricane Brenda :

\begin{verbatim}
;BRENDA   *092345z4903.50N\07202.75W@088/036/HC/150^200/0980>090&030%040
;BRENDA   *100045z4905.50N/07201.75W@101/047/HC/104^123/0980>065&020%040
\end{verbatim}

\section{Bulletins NWS}\label{bulletins-nws}

APRS supporte la diffusion des bulletins National Weather Service. Voir
chapitre 14.

